%% kapa: exam=1 round=20 year=2009 period=2 form=A subject=accounting subject_ko=회계학 seq=18 number=18
%% source: Q-Net 감정평가사 기출문제 (cst00309) artlSeq=5024178 file=extracted/1008/5024178/2009년 제20회 감정평가사 1차시험 기출문제/2교시.hwp
%% conversion: kordoc 4.24.51 -> tools/gen1.py (auto); review=proofread-sonnet; flags=none
\begin{question}{18}
(주)한국은 재고자산에 대해 실지재고조사법을 적용하고 있다. 20×9년 7월중 상품거래자료가 다음과 같을 때 7월의 선입선출법과 평균법에 의한 매출원가 차액은 얼마인가? (단, 장부상의 재고와 실지재고는 일치하며, 재고자산평가손실은 없다.)

{\small\centering
\begin{tabular}{cc r r r}
\hline
일자 & 적요 & \multicolumn{1}{c}{수량} & \multicolumn{1}{c}{매입단가} & \multicolumn{1}{c}{금액} \\
\hline
7월 1일 & 월초재고 & 200개 & ￦100 & ￦20,000 \\
7월 10일 & 매\quad 입 & 300 & 110 & 33,000 \\
7월 15일 & 매\quad 출 & 400 & - & - \\
7월 20일 & 매\quad 입 & 200 & 120 & 24,000 \\
7월 25일 & 매\quad 출 & 200 & - & - \\
\hline
\end{tabular}\par}
\begin{choices}
  \item ￦0
  \item ￦500
  \item ￦1,000
  \item ￦1,500
  \item ￦2,000
\end{choices}
\end{question}
\answer{3}
